\documentclass[12pt]{article}

\usepackage[utf8]{inputenc}
\usepackage[T1]{fontenc}
\usepackage[polish,provide=*]{babel}
\usepackage{lmodern}
\usepackage{amsmath}
\usepackage{latexsym,amsfonts,amssymb,amsthm,amsmath}
\usepackage{enumitem}
\usepackage{hyperref}
\usepackage{float}
\usepackage{graphicx}
\usepackage{subcaption}
\usepackage{booktabs}
\usepackage{physics}
\graphicspath{{./images/}}

\setlength{\parindent}{0in}
\setlength{\oddsidemargin}{0in}
\setlength{\textwidth}{6.5in}
\setlength{\textheight}{8.8in}
\setlength{\topmargin}{0in}
\setlength{\headheight}{18pt}

\title{Analysis of open transmon system}
\author{Kacper Kłos}

\begin{document}

\maketitle

\section{Transmon and harmonic oscillator}

This section will derive and analyse behaviour of simple system consisting of a transmon modeled as Josephson junction and quantum harmonic oscillator.

\subsection{The Josephson junction}

The Josephson junction is made with two disconnected superconductors between which there is a layer of oxide. This setup allows Cooper pairs to tunnel throughout the junctin. 

To analyse the behaviour of Josephson junction we will only need two equations. One for the junction energy:
\begin{equation}
    E = E_J \cos{\varphi}
\end{equation}
And for the voltage across the junction:
\begin{equation}
    V = \frac{\Phi_0}{2\pi} \dv{\varphi}{t}
\end{equation}

\subsection{Quantum LC circuit}

We have a classical LC circuit with the Hamiltonian:
\begin{equation}
    H_{LC} = \frac{Q^2}{2 C} + \frac{\Phi^2}{2 L}
\end{equation}







\end{document}
